%% kapa: exam=1 round=23 year=2012 period=2 form=A subject=appraisal_law subject_ko=부동산관계법규 seq=35 number=75
%% source: Q-Net 감정평가사 기출문제 (cst00309) artlSeq=5052580 file=raw/1008/5052580/2012년 제23회 감정평가사 제1차시험 기출문제(2교시 A형).pdf
%% conversion: kordoc 4.24.51 -> tools/gen1.py (auto); review=proofread-sonnet; flags=none
\begin{question}{75}
측량ㆍ수로조사 및 지적에 관한 법령상 지목에 관한 설명으로 옳지 않은 것은?
\begin{choices}
  \item 묘지의 관리를 위한 건축물 부지의 지목은 “대”로 한다.
  \item 과수류를 집단적으로 재배하는 토지에 접속된 주거용 건축물 부지의 지목은 “대”로 한다.
  \item 「주차장법」에 따른 노상주차장 부지의 지목은 “주차장”으로 한다.
  \item 제조업을 하고 있는 공장시설물의 부지와 같은 구역에 있는 부속시설물인 의료시설 부지의 지목은 “공장용지”로 한다.
  \item 자연의 유수가 있을 것으로 예상되는 토지의 지목은 “하천”으로 한다.
\end{choices}
\end{question}
\answer{3}
